% =====================================================================
%  Chapitre 5 — Sprint 3 : Assistant d'intelligence artificielle agentique
% =====================================================================
\chapter{Sprint~3 : Assistant d'intelligence artificielle agentique}
\label{chap:sprint3}

\section*{Introduction}
Ce sprint porte la principale valeur ajoutée du projet : un assistant
conversationnel capable de générer et de modifier des modèles d'e-mail en
manipulant de véritables blocs éditables, et même de reconstruire une
campagne à partir d'une affiche. Nous présentons d'abord les objectifs du
sprint, puis la conception de l'assistant — modèles employés, architecture
de génération, protocole d'exposition des outils et chaîne de traitement de
l'image —, avant de décrire les diagrammes de séquence et la réalisation.

\section{Objectifs du sprint}
Ce sprint couvre les epics~7 à~10 du backlog. Le
tableau~\ref{tab:backlog-s3} en détaille les user stories.

\begin{table}[H]
\centering
\caption{Backlog du sprint~3}
\label{tab:backlog-s3}
\small
\begin{tabularx}{\linewidth}{clX}
\toprule
\textbf{ID} & \textbf{User story} & \textbf{Description} \\
\midrule
7.1 & Générer par consigne & En tant qu'éditeur, je veux qu'un e-mail complet et à ma marque soit généré à partir d'une consigne en langage naturel. \\
8.1 & Éditer en conversation & En tant qu'éditeur, je veux modifier mon e-mail par dialogue (thème, image, disposition, ajout ou déplacement de blocs). \\
8.2 & Éditer par sélection & En tant qu'éditeur, je veux sélectionner un élément et lui appliquer une consigne ciblée. \\
9.1 & Affiche vers e-mail & En tant qu'éditeur, je veux importer une affiche et obtenir un e-mail éditable reprenant la campagne. \\
10.1 & Exposition des outils (MCP) & En tant que système, je veux exposer les outils de construction via un protocole standardisé afin que le modèle les découvre dynamiquement. \\
10.2 & Modèle auto-hébergé & En tant qu'entreprise, je veux que l'inférence s'exécute sur un serveur privé afin de préserver la confidentialité des données. \\
\bottomrule
\end{tabularx}
\end{table}

\section{Conception de l'assistant}

\subsection{Le modèle de blocs comme contrat d'outils}
Le principe fondateur de l'assistant est le suivant : au lieu de produire un
bloc de HTML monolithique, le modèle de langage \textbf{appelle des outils}
de construction qui renvoient exactement le même format de blocs
(\texttt{BlockData}) que l'éditeur manuel décrit au
chapitre~\ref{chap:sprint2}. Le HTML n'est généré qu'au moment de l'export.
Cette approche, dite \emph{agentique}, présente deux avantages décisifs :
les modifications visuelles et les modifications par IA opèrent sur la même
structure de données (elles sont donc interchangeables), et le résultat est
immédiatement et finement éditable par l'utilisateur.

Un ensemble de \textbf{32 outils} a été défini : un outil par type de bloc,
complété par des outils transversaux (définition du thème, ouverture de
section, carte, en-tête sur image, barre de couleurs de marque, changement
de disposition, déplacement, édition et suppression de blocs).

\subsection{Modèles et infrastructure}
L'inférence repose sur un \textbf{serveur de modèle de langage
auto-hébergé} exploitant le moteur haute performance vLLM, exposé via une
API compatible OpenAI et protégé par authentification. Deux modèles à poids
ouverts y sont servis : \textbf{Gemma} (modèle par défaut) et un modèle de
repli, tous deux dotés d'un large contexte. Ce choix d'auto-hébergement est
central : il garantit qu'aucune donnée de l'entreprise n'est transmise à une
API tierce, répondant ainsi à l'exigence de souveraineté formulée dans les
besoins non fonctionnels.

\subsection{Architecture de génération : planificateur et exécuteur}
Le chemin principal de génération d'un e-mail neuf repose sur un
\textbf{planificateur déterministe}. Plutôt que de laisser le modèle
enchaîner librement des dizaines d'appels d'outils — approche fragile avec
un modèle de taille modérée —, une \textbf{unique requête contrainte par un
schéma} produit une spécification de design (\emph{DesignSpec}) : système de
design, ambiance, palette de couleurs et liste ordonnée de sections. Cette
spécification est ensuite \textbf{exécutée par le code} pour construire les
blocs de manière déterministe. Il en résulte une génération rapide, sans
duplication et avec un pied de page unique. Lorsque le planificateur échoue
ou renvoie un résultat trop pauvre, une \textbf{voie agentique de repli}
prend le relais, laissant le modèle appeler directement les outils. La
figure~\ref{fig:archi-ia} synthétise cette architecture.

% [FIGURE : architecture de l'assistant IA — source architecture-ia.drawio]
\figplaceholder{Architecture de l'assistant d'intelligence artificielle}{fig:archi-ia}{7.5cm}

\subsection{Exposition des outils via le protocole MCP}
Pour les tours pilotés par le modèle (édition et repli agentique), les
outils de construction sont exposés au moyen du \textbf{Model Context
Protocol} (MCP). Un point d'accès dédié agit comme serveur MCP : le modèle y
\textbf{découvre dynamiquement} les 32 outils disponibles plutôt que de se
les voir énumérer dans l'invite. Cette décorrélation entre la définition des
outils et l'invite constitue une architecture agentique moderne et
standardisée. Le serveur MCP réutilise exactement les mêmes définitions
d'outils que l'éditeur, garantissant une source de vérité unique.

\subsection{Chaîne « affiche vers e-mail »}
La fonction d'import d'affiche traite une image en \textbf{trois étapes}.
\begin{enumerate}[leftmargin=1.4cm]
  \item \textbf{Lecture par vision.} L'image est analysée par le modèle
  multimodal, qui en extrait un rapport structuré : texte exact (OCR),
  offre (prix, cadeaux, forfaits, code promotionnel, appel à l'action),
  couleurs dominantes, ambiance, marque et disposition. Cette étape, la seule
  réellement multimodale, s'appuie sur un décodage contraint par schéma pour
  fiabiliser l'extraction.
  \item \textbf{Rapport vers directive.} Le rapport et la consigne de
  l'utilisateur sont transformés en une directive de génération et en une
  amorce de thème, qui alimentent le générateur de blocs déjà décrit :
  l'image gouverne la marque et le contenu, la consigne oriente la
  structure.
  \item \textbf{Passe critique.} Une passe d'auto-vérification compare
  l'e-mail produit à l'affiche d'origine et signale les écarts (prix erroné,
  élément manquant), corrigés par une passe finale.
\end{enumerate}
Il est essentiel de noter que la plateforme ne colle jamais l'affiche : elle
\textbf{reconstruit} nativement la campagne en blocs éditables. La
figure~\ref{fig:pipeline-image} illustre cette chaîne de traitement.

% [FIGURE : chaîne de traitement « affiche → e-mail »]
\figplaceholder{Chaîne de traitement « affiche vers e-mail »}{fig:pipeline-image}{6cm}

\subsection{Intelligence de design et fiabilité}
Plusieurs mécanismes garantissent la qualité et la fiabilité du résultat.
\begin{itemize}[leftmargin=1.4cm]
  \item \textbf{Cinq systèmes de design} (éditorial, audacieux, minimal,
  luxe, corporate) paramètrent l'échelle typographique, les alignements, les
  espacements et le traitement des boutons et des cartes~; le système est
  choisi en fonction du sujet.
  \item Un \textbf{dériveur de palette} et une \textbf{garde de contraste}
  assurent que le texte reste toujours lisible sur son fond, quelle que soit
  la couleur choisie.
  \item Un \textbf{routeur d'intention} classe chaque tour (effacer,
  réécrire, image, éditer, changer de thème) afin d'aiguiller la requête vers
  le traitement — déterministe ou piloté par le modèle — le plus adapté.
  \item Des \textbf{confirmations honnêtes} garantissent que l'assistant
  n'annonce un succès que si le modèle a réellement été modifié, un
  compteur de mutations servant de preuve de changement.
\end{itemize}

\section{Diagrammes de séquence}

\subsection{Séquence « Générer un e-mail par consigne »}
La figure~\ref{fig:seq-generate} décrit le flux de génération : la consigne
est transmise à la route d'IA, le planificateur produit une spécification de
design via le serveur d'inférence, le code exécute cette spécification en
blocs, les images sont résolues, puis le modèle est renvoyé au client et
appliqué au canevas.

% [FIGURE : diagramme de séquence « Générer un e-mail par consigne »]
\figplaceholder{Diagramme de séquence « Générer un e-mail par consigne »}{fig:seq-generate}{7.5cm}

\subsection{Séquence « Reconstruire depuis une affiche »}
La figure~\ref{fig:seq-image} illustre la chaîne complète de l'import
d'affiche : lecture par vision, transformation en directive, génération des
blocs et passe critique.

% [FIGURE : diagramme de séquence « Reconstruire depuis une affiche »]
\figplaceholder{Diagramme de séquence « Reconstruire depuis une affiche »}{fig:seq-image}{7.5cm}

\section{Réalisation}

\subsection{Panneau de discussion et génération}
L'assistant se présente sous la forme d'un panneau de discussion intégré à
l'éditeur d'e-mail (figure~\ref{fig:ui-ai-chat}). À partir d'une simple
consigne — par exemple « crée un e-mail Black Friday pour notre boutique » —
l'assistant produit en une dizaine de secondes un e-mail complet et à la
marque : en-tête, section d'offre, grille de produits et pied de page.

% [FIGURE : capture — panneau de discussion IA et e-mail généré]
\figplaceholder{Panneau de discussion IA et e-mail généré}{fig:ui-ai-chat}{6.5cm}

\subsection{Édition conversationnelle et ciblée}
L'utilisateur peut ensuite affiner son e-mail par le dialogue : « passe-le
en thème sombre », « change la première image », « déplace la section des
produits avant l'offre », « ajoute une boîte de prix ». L'assistant vise le
bon bloc et applique la modification. Il est également possible de
\textbf{sélectionner} un bloc ou une section dans le canevas puis de saisir
une consigne : la modification s'applique alors précisément à l'élément
sélectionné. La figure~\ref{fig:ui-ai-edit} illustre l'édition ciblée par
sélection.

% [FIGURE : capture — édition ciblée par sélection]
\figplaceholder{Édition ciblée par sélection}{fig:ui-ai-edit}{6cm}

\subsection{Import d'une affiche}
La figure~\ref{fig:ui-image-import} présente l'import d'une affiche : après
téléversement, l'assistant affiche l'avancement des étapes (lecture, analyse,
génération) puis produit l'e-mail éditable reconstruit, fidèle à la marque
et à l'offre de l'affiche d'origine.

% [FIGURE : capture — import d'affiche et e-mail reconstruit]
\figplaceholder{Import d'une affiche et e-mail reconstruit}{fig:ui-image-import}{6.5cm}

\section*{Conclusion}
Ce sprint a doté la plateforme de son principal différenciateur : un
assistant agentique qui génère et modifie de véritables blocs éditables,
reconstruit une campagne à partir d'une affiche grâce à la vision, et
s'appuie sur un modèle de langage auto-hébergé garantissant la souveraineté
des données. L'architecture combine un planificateur déterministe rapide,
une voie agentique de repli et une exposition des outils via le protocole
MCP, le tout renforcé par une intelligence de design et des mécanismes de
fiabilité. Le sprint suivant transforme cette plateforme fonctionnellement
riche en un service commercialisable : abonnements, intégrations et mise en
production.
